\documentclass{standalone}
\usepackage{circuitikz}
\usetikzlibrary{calc}
\definecolor{circuitnotemark}{HTML}{FE00FE}
\begin{document}
\begin{circuitikz}[american, line width=0.8pt]
\ctikzset{bipoles/length=1.2cm}
\foreach \x in {0,2,4} {\foreach \y in {0,-2,-4,-6} {\fill[gray, opacity=0.35] (\x,\y) circle (1.1pt);}}
\node[gray, font=\scriptsize] at (0,0.84) {1};
\node[gray, font=\scriptsize] at (2,0.84) {2};
\node[gray, font=\scriptsize] at (4,0.84) {3};
\node[gray, font=\scriptsize, anchor=east] at (-0.84,0) {a};
\node[gray, font=\scriptsize, anchor=east] at (-0.84,-2) {b};
\node[gray, font=\scriptsize, anchor=east] at (-0.84,-4) {c};
\node[gray, font=\scriptsize, anchor=east] at (-0.84,-6) {d};
\coordinate (b1) at (0,-2);
\coordinate (d1) at (0,-6);
\coordinate (b2) at (2,-2);
\coordinate (b3) at (4,-2);
\coordinate (d3) at (4,-6);
\draw (b1) to[esource, l_=$\dot{E}$] (d1); % line 3
\draw (-0.18,-4) sin ++(0.09,0.09) cos ++(0.09,-0.09) sin ++(0.09,-0.09) cos ++(0.09,0.09); % line 3
\draw (b1) to[nos, l_=$\mathrm{SW}$] (b2); % line 4
\draw (b2) to[R, l_=$\dot{Z}_{L}$] (b3); % line 5
\draw (b3) to[R, l_=$R_{L}$] (d3); % line 6
\draw (d1) -- (d3); % line 8
\node[anchor=south west, inner sep=2pt, circuitnotemark, font=\LARGE\bfseries] (circuittitle) at (current bounding box.north west) {X};
\path (circuittitle.south west) rectangle ++(10.455,0.853);
\end{circuitikz}
\end{document}
